% RESUMO
O Aprendizado Federado permite treinar modelos em dispositivos de borda sem centralizar dados brutos, mas a heterogeneidade estatística entre clientes (dados não-IID) degrada a convergência e a acurácia do algoritmo padrão, o FedAvg. O FedProx propõe mitigar o problema com um termo proximal na função objetivo local, porém a literatura recente apresenta resultados contraditórios sobre quando ele de fato supera o FedAvg e raramente mede o custo computacional que impõe aos clientes. Este trabalho é uma pesquisa experimental que compara FedAvg e FedProx, tendo FedAvgM e FedAdam como referências (e SCAFFOLD, se integrado a tempo), em simulações com o framework Flower sobre CIFAR-10 e Fashion-MNIST. A heterogeneidade é controlada por partição de Dirichlet ($\alpha \in \{0{,}05; 0{,}1; 0{,}5; 1; 100\}$), partição patológica e desbalanceamento de quantidade, e o hiperparâmetro proximal variará em $\mu \in \{0{,}01; 0{,}1; 1\}$, com sementes pareadas (cinco exploratórias e dez confirmatórias), teste de Wilcoxon e teste de equivalência. Serão avaliados acurácia global, estabilidade, rodadas até a acurácia-alvo, dispersão entre clientes e tempo de computação local, este último também em uma CPU restrita. A contribuição esperada é um mapa de trade-off $\alpha \times \mu \times$ (acurácia, custo) que indica em quais níveis e tipos de heterogeneidade o FedProx compensa. A proposta é um protocolo de avaliação pareada em que, na mesma semente, os algoritmos compartilham partição, inicialização e clientes sorteados, de modo que o FedProx com $\mu = 0$ reproduz o FedAvg de forma exata. Nas 285 execuções já realizadas no Fashion-MNIST, o FedProx com o $\mu^* = 0{,}01$ escolhido na fase exploratória ficou, na média de cada nível de heterogeneidade, a no máximo 0,02~p.p.\ do FedAvg e, na fase confirmatória, mostrou-se estatisticamente equivalente a ele nos cinco níveis da partição de Dirichlet, na partição patológica e sob desvio de quantidade, sem o ganho previsto para a heterogeneidade forte, enquanto o FedAdam, uma correção no servidor, melhorou a acurácia em 1,5 a 5,2~p.p.

Repositório do projeto: \linkProjeto
